\apendice{Plan de Proyecto Software}

\section{Introducción}

\section{Planificación temporal}

\subsection{Planificación por Sprints}

Siguiendo la metodología ágil de Scrum, la planificación temporal a lo largo de este proyecto va a ser realizada mediante sprints.

\begin{Sprint}{1}
	
\bloque{Duración}
La duración definida para este primer sprint fue de dos semanas, que inició el 16 de septiembre de 2025 y finalizó el 29 de septiembre de 2025.

\bloque{Objetivos}
Durante la reunión de planificación de este sprint se propusieron los siguientes objetivos:
\begin{enumerate}
	\item Comenzar a realizar la base de datos.
	\item Descargar los programas necesarios para la documentación del proyecto en LaTeX.
	\item Investigar y revisar apuntes sobre la metodología Scrum.
	\item Investigar y revisar apuntes relacionado a los casos de uso.
\end{enumerate}

\bloque{Observaciones}
El comienzo de la base de datos se realizó desde un boceto de diagrama entidad-relación realizado durante la reunión.

\bloque{Revisión del Sprint}
Se lograron realizar todos los objetivos propuestos en la reunión de planificación, aunque queda profundizar más sobre la recopilación de información respecto al funcionamiento de los casos de uso.
\end{Sprint}

\begin{Sprint}{2}
	
\bloque{Duración}
A este sprint, se le ha especificado una duración de una semana, que va desde el 30 de septiembre de 2025 hasta el 6 de cctubre de 2025.
	
\bloque{Objetivos}
Durante la reunión de planificación de este sprint se propusieron los siguientes objetivos:
\begin{enumerate}
	\item Realizar el .gitignore
	\item Subir la plantilla de LaTeX al repositorio
	\item Hacer el diagrama entidad-relación
	\item Realizar el modelo relacional
	\item Comenzar con el listado de casos de uso
\end{enumerate}
	
%\bloque{Observaciones}
	
\bloque{Revisión del Sprint}
En este sprint, se han conseguido todos los objetivos propuesto al comienzo del mismo durante la reunión de planificación, aunque hay que realizar mejoras con el listado de casos de uso de la aplicación.
\end{Sprint}

\begin{Sprint}{3}
	
\bloque{Duración}
Al igual que el sprint anterior, este tercer sprint tiene una duación de una semana. Comenzando el 7 de octubre de 2025 y finalizando el 13 de octubre de 2025.

\bloque{Objetivos}
Durante la reunión de planificación de este sprint se propusieron los siguientes objetivos:
\begin{enumerate}
	\item Definir lenguaje de programación y framework para el proyecto
	\item Definir arquitectura 
	\item Documentar los tres primeros sprints
	\item Definir casos de uso a alto nivel
	\item Definir y extender un caso de uso
\end{enumerate}

\bloque{Revisión del Sprint}

\end{Sprint}

\section{Estudio de viabilidad}

\subsection{Viabilidad económica}

\subsection{Viabilidad legal}


